% Generated by tools/make_body.py; edit generator, not this file.
% UNPROOFREAD OCR TRANSCRIPTION: page-image verification remains required.
\chapter{Elliptic Functions}\label{chap:7}

% Original PDF physical page 279
\sourcepage{279}
\setcounter{section}{0}\section{Simply Periodic Functions}\index{Elliptic function|(}\label{sec:7.1}
A function $f(z)$ is said to be periodic with period\index{Period} $\omega \neq 0$ if
\begin{equation*}
f (z + \omega) = f (z)
\end{equation*}
for all z. For instance, $e^{z}$ has the period $2\pi i$, and $\sin z$ and $\cos z$ have the period $2\pi$. To be more precise, we are interested only in analytic or meromorphic functions $f(z)$, and they shall be considered in a region $\Omega$ which is mapped onto itself by the translation $z \to z + \omega$.

If $\omega$ is a period, so are all integral multiples $n\omega$. There may be other periods as well, but for the present we focus our attention exclusively on the periods $n\omega$. From this point of view we shall call $f(z)$ a \emph{simply periodic function} with period $\omega$. In particular, it is irrelevant whether $\omega$ is itself a multiple of another period.

\setcounter{section}{1}\setcounter{subsection}{0}\subsection{Representation by Exponentials}\label{sec:7.1.1}
The simplest function with period $\omega$ is the exponential $e^{2\pi iz/\omega}$. It is a fundamental fact that any function with period $\omega$ can be expressed in terms of this particular function.

Let $\Omega$ be a region with the property that $z \in \Omega$ implies $z + \omega \in \Omega$ and $z - \omega \in \Omega$. We define $\Omega'$ in the $\zeta$-plane to be the image of $\Omega$ under the mapping $\zeta = e^{2\pi iz / \omega}$; it is obviously a region. For instance, if $\Omega$ is the whole plane, then $\Omega'$ is the plane punctured at 0. If $\Omega$ is a parallel strip, defined by $a < \operatorname{Im}(2\pi z / \omega) < b$, then $\Omega'$ is the annulus $e^{-b} < |\zeta| < e^{-a}$.


% Original PDF physical page 280
\sourcepage{280}
Suppose that $f(z)$ is meromorphic in $\Omega$ and has the period $\omega$. Then there exists a unique function $F$ in $\Omega'$ such that
\begin{equation*}
f (z) = F \left(e ^ {2 \pi i z / \omega}\right).\tag{1}\label{eq:7.1}
\end{equation*}
Indeed, to determine $F(\zeta)$ we write $\zeta = e^{2\pi iz / \omega}$; $z$ is unique up to an additive multiple of $\omega$, and this multiple does not influence the value $f(z)$. It is evident that $F$ is meromorphic. Conversely, if $F$ is meromorphic in $\Omega'$, then \hyperref[eq:7.1]{(1)} defines a meromorphic function $f$ with period $\omega$.

\setcounter{section}{1}\setcounter{subsection}{1}\subsection{The Fourier Development}\label{sec:7.1.2}
Assume that $\Omega'$ contains an annulus $r_1 < |\zeta| < r_2$ in which $F$ has no poles. In this annulus $F$ has a Laurent development
\begin{equation*}
F (\zeta) = \sum_ {n = - \infty} ^ {\infty} c _ {n} \zeta^ {n},
\end{equation*}
and we obtain
\begin{equation*}
f (z) = \sum_ {- \infty} ^ {\infty} c _ {n} e ^ {2 \pi i n z / \omega}.
\end{equation*}
This is the complex Fourier development\index{Fourier development} of $f(z)$, valid in the parallel strip that corresponds to the given annulus.

The coefficients (cf. Chap.~\hyperref[chap:5]{5}, Sec.~\hyperref[sec:5.1.3]{1.3}) are given by
\begin{equation*}
c _ {n} = \frac {1}{2 \pi i} \int_ {| \zeta | = r} F (\zeta) \zeta^ {- n - 1} d \zeta , \quad (r _ {1} <   r <   r _ {2}),
\end{equation*}
and by change of variable this becomes
\begin{equation*}
c _ {n} = \frac {1}{\omega} \int_ {a} ^ {a + \omega} f (z) e ^ {- 2 \pi i n z / \omega} d z.
\end{equation*}
Here a is an arbitrary point in the parallel strip, and the integration is along any path from a to $a + \omega$ which remains within the strip. If $f(z)$ is analytic in the whole plane, the same Fourier development is valid everywhere.

\setcounter{section}{1}\setcounter{subsection}{2}\subsection{Functions of Finite Order}\label{sec:7.1.3}
When $\Omega$ is the whole plane $F(\zeta)$ has isolated singularities at $\zeta = 0$ and $\zeta = \infty$. If both these singularities are inessential, that is, either removable singularities or poles, then F is a rational function. We say in this case that f has finite order, equal to the order of F.

We recall that a rational function assumes every complex value, including $\infty$, the same number of times, provided that we observe the usual multiplicity convention. We obtain a similar result for simply periodic functions of finite order if we agree not to distinguish between z
% Original PDF physical page 281
\sourcepage{281}
and\index{Discrete set|(}\index{Set!discrete|(} $z + \omega$. For convenient terminology, let us say that $z + n\omega$ is equivalent to z. If f is of order m we find that every complex value $c \neq F(0)$ and $F(\infty)$ is assumed at m inequivalent points, with due count of multiplicities. We observe further that $f(z) \to F(0)$ when $\operatorname{Im}(z/\omega) \to -\infty$ and $f(z) \to F(\infty)$ when $\operatorname{Im}(z/\omega) \to \infty$. If we are willing to agree that these values are also “assumed” (with proper multiplicity), we can maintain that all complex values are assumed exactly m times.

For another interpretation we may consider the period strip, defined by $0 \leq \operatorname{Im}(z/\omega) < 2\pi$. Since this strip contains only one representative from each equivalence class we find that $f(z)$ assumes each complex value m times in the period strip, except that the values $F(0)$ and $F(\infty)$ require a special convention.

\setcounter{section}{1}\section{Doubly Periodic Functions}\label{sec:7.2}
The terms elliptic function\index{Elliptic function} and doubly periodic function\index{Doubly periodic function} are interchangeable; we have already met examples of such functions in connection with the conformal mapping of rectangles and certain triangles (Chap.~\hyperref[chap:6]{6}, Sec.~\hyperref[sec:6.2]{2}). Elliptic functions have been the object of very extensive study, partly because of their function theoretic properties and partly because of their importance in algebra and number theory. Our introduction to the topic covers only the most elementary aspects.

\setcounter{section}{2}\setcounter{subsection}{0}\subsection{The Period Module}\label{sec:7.2.1}
Let $f(z)$ be meromorphic in the whole plane. We shall examine the set M of all its periods. If $\omega$ is a period, so are all integral multiples $n\omega$, and if $\omega_{1}$ and $\omega_{2}$ belong to M, so does $\omega_{1} + \omega_{2}$; as a consequence, all linear combinations $n_{1}\omega_{1} + n_{2}\omega_{2}$ are in M. In algebra, a set with these properties is called a module\index{Module} (more precisely: a module over the integers), and we shall call M the period module of f.

Apart from the trivial case of a constant function, M has also a topological property: all its points are isolated. In fact, since $f(\omega) = f(0)$ for all $\omega \in M$ the existence of a finite accumulation point would immediately imply that f is constant. A module with isolated points is said to be discrete.

Our first step is to determine all discrete modules.

\begin{originalstatement}{Theorem}{1}\label{thm:7.1}
A discrete module consists either of zero alone, of the integral multiples $n\omega$ of a single complex number $\omega \neq 0$, or of all linear combinations $n_1\omega + n_2\omega_2$ with integral coefficients of two numbers $\omega_1, \omega_2$ with nonreal ratio $\omega_2 / \omega_1$.
\end{originalstatement}
As soon as M contains a number $\omega \neq 0$ it also contains one, call it $\omega_{1}$, whose absolute value is a minimum. Indeed, if r is large enough the
% Original PDF physical page 282
\index{Discrete set|)}\index{Set!discrete|)}\sourcepage{282}
disk $|z| \leq r$ contains a point from M, other than 0. Because the points are isolated there are only a finite number of such points, and we choose $\omega_{1}$ to be one closest to the origin (the reader may show that there are always 2, 4, or 6 closest points). The multiples $n\omega_{1}$ are also in M, and these may be all.

Suppose now that there exists an $\omega \in M$ which is not an integral multiple of $\omega_{1}$. Among all such there is one, $\omega_{2}$, whose absolute value is smallest. We claim that $\omega_{2}/\omega_{1}$ is not real. If it were, there would exist an integer n such that $n < \omega_{2}/\omega_{1} < n + 1$. This would give $0 < |n\omega_{1} - \omega_{2}| < |\omega_{1}|$, an obvious contradiction.

It can now be concluded that all numbers in M are of the form $n_{1}\omega_{1} + n_{2}\omega_{2}$. First of all, because $\omega_{2}/\omega_{1}$ is nonreal any complex number $\omega$ can be written in the form $\lambda_{1}\omega_{1} + \lambda_{2}\omega_{2}$ with real $\lambda_{1}$ and $\lambda_{2}$. To see this we need only solve the equations
\begin{equation*}
\begin{array}{l} \omega = \lambda_ {1} \omega_ {1} + \lambda_ {2} \omega_ {2} \\ \bar {\omega} = \lambda_ {1} \bar {\omega} _ {1} + \lambda_ {2} \bar {\omega} _ {2}. \end{array}
\end{equation*}
Since the determinant $\omega_1\bar{\omega}_2 - \omega_2\bar{\omega}_1$ is $\neq 0$ the system has a unique solution $(\lambda_1,\lambda_2)$; but $(\bar{\lambda}_1,\bar{\lambda}_2)$ is also a solution, and we conclude that $\lambda_1$ and $\lambda_2$ are real. To continue the proof, there exist integers $m_1, m_2$ such that $|\lambda_1 - m_1| \leq \frac{1}{2}, |\lambda_2 - m_2| \leq \frac{1}{2}$. If $\omega$ belongs to $M$, so does
\begin{equation*}
\omega^ {\prime} = \omega - m _ {1} \omega_ {1} - m _ {2} \omega_ {2}.
\end{equation*}
We have $|\omega^{\prime}| < \frac{1}{2} |\omega_{1}| + \frac{1}{2} |\omega_{2}| \leq |\omega_{2}|$ where the first inequality is strict because $\omega_{2}$ is not a real multiple of $\omega_{1}$. By the way $\omega_{2}$ was chosen it follows that $\omega^{\prime}$ must be an integral multiple of $\omega_{1}$, and hence $\omega$ has the asserted form.

\setcounter{section}{2}\setcounter{subsection}{1}\subsection{Unimodular Transformations}\index{Unimodular transformation|(}\index{Unimodular transformation}\label{sec:7.2.2}
We assume henceforth that it is the third alternative in Theorem~\hyperref[thm:7.1]{1} that occurs. The pair $(\omega_{1},\omega_{2})$ has the property that any $\omega \in M$ has a unique representation of the form $\omega = n_1\omega_1 + n_2\omega_2$. Any pair with this property will be called a basis of $M$ (even if it is not determined by the construction in the proof of Theorem~\hyperref[thm:7.1]{1}).

We investigate the relation between two bases $(\omega_{1},\omega_{2})$ and $(\omega_1',\omega_2')$. Because $(\omega_{1},\omega_{2})$ is a basis there exist integers $a, b, c, d$ such that
\begin{equation*}
\begin{array}{r} \omega_ {2} ^ {\prime} = a \omega_ {2} + b \omega_ {1} \\ \omega_ {1} ^ {\prime} = c \omega_ {2} + d \omega_ {1}. \end{array}\tag{2}\label{eq:7.2}
\end{equation*}
We prefer to write these equations in matrix form
\begin{equation*}
\binom{\omega_ {2} ^ {\prime}}{\omega_ {1} ^ {\prime}} = \left( \begin{array}{c c} a & b \\ c & d \end{array} \right) \binom{\omega_ {2}}{\omega_ {1}}.
\end{equation*}
% Original PDF physical page 283
\sourcepage{283}
The same relation is valid for the complex conjugates, and we have thus
\begin{equation*}
\left( \begin{array}{c c} \omega_ {2} ^ {\prime} & \bar {\omega} _ {2} ^ {\prime} \\ \omega_ {1} ^ {\prime} & \bar {\omega} _ {1} ^ {\prime} \end{array} \right) = \left( \begin{array}{c c} a & b \\ c & d \end{array} \right) \left( \begin{array}{c c} \omega_ {2} & \bar {\omega} _ {2} \\ \omega_ {1} & \bar {\omega} _ {1} \end{array} \right).\tag{3}\label{eq:7.3}
\end{equation*}
Since $(\omega_1', \omega_2')$ is also a basis we have similarly
\begin{equation*}
\left( \begin{array}{c c} \omega_ {2} & \bar {\omega} _ {2} \\ \omega_ {1} & \bar {\omega} _ {1} \end{array} \right) = \left( \begin{array}{c c} a ^ {\prime} & b ^ {\prime} \\ c ^ {\prime} & d ^ {\prime} \end{array} \right) \left( \begin{array}{c c} \omega_ {2} ^ {\prime} & \bar {\omega} _ {2} ^ {\prime} \\ \omega_ {2} ^ {\prime} & \bar {\omega} _ {1} ^ {\prime} \end{array} \right)\tag{4}\label{eq:7.4}
\end{equation*}
with integral $a', b', c', d'$.

From \hyperref[eq:7.3]{(3)} and \hyperref[eq:7.4]{(4)} we obtain
\begin{equation*}
\left( \begin{array}{c c} \omega_ {2} & \bar {\omega} _ {2} \\ \omega_ {1} & \bar {\omega} _ {1} \end{array} \right) = \left( \begin{array}{c c} a ^ {\prime} & b ^ {\prime} \\ c ^ {\prime} & d ^ {\prime} \end{array} \right) \left( \begin{array}{c c} a & b \\ c & d \end{array} \right) \left( \begin{array}{c c} \omega_ {2} & \bar {\omega} _ {2} \\ \omega_ {1} & \bar {\omega} _ {1} \end{array} \right).\tag{5}\label{eq:7.5}
\end{equation*}
Here the determinant $\omega_{2}\bar{\omega}_{1} - \omega_{1}\bar{\omega}_{2}$ is $\neq 0$, for otherwise any two numbers in the module would have a real ratio, contrary to assumption. A matrix with determinant $\neq 0$ has an inverse matrix, and if we multiply \hyperref[eq:7.5]{(5)} by the inverse of $\left( \begin{array}{cc}\omega_2 & \bar{\omega}_2\\ \omega_1 & \bar{\omega}_1 \end{array} \right)$ we obtain
\begin{equation*}
\left( \begin{array}{c c} a ^ {\prime} & b ^ {\prime} \\ c ^ {\prime} & d ^ {\prime} \end{array} \right) \left( \begin{array}{c c} a & b \\ c & d \end{array} \right) = \left( \begin{array}{c c} 1 & 0 \\ 0 & 1 \end{array} \right).
\end{equation*}
The matrices $\begin{pmatrix} a & b \\ c & d \end{pmatrix}$ and $\begin{pmatrix} a' & b' \\ c' & d' \end{pmatrix}$ are inverse to each other. In particular, their determinants must satisfy
\begin{equation*}
\left| \begin{array}{c c} a ^ {\prime} & b ^ {\prime} \\ c ^ {\prime} & d ^ {\prime} \end{array} \right| \cdot \left| \begin{array}{c c} a & b \\ c & d \end{array} \right| = 1,
\end{equation*}
and since both are integers we must have
\begin{equation*}
\left| \begin{array}{c c} a & b \\ c & d \end{array} \right| = \left| \begin{array}{c c} a ^ {\prime} & b ^ {\prime} \\ c ^ {\prime} & d ^ {\prime} \end{array} \right| = \pm 1.
\end{equation*}
Linear transformations of the form \hyperref[eq:7.2]{(2)} with integral coefficients and determinant $\pm1$ are said to be unimodular. We have proved:

\begin{originalassertion}
Any two bases of the same module are connected by a unimodular transformation\index{Unimodular transformation}.
\end{originalassertion}

Geometrically, it is natural to consider the parallelogram spanned by a basis $(\omega_{1},\omega_{2})$ in its relation to the lattice formed by all numbers in the module. Figure~\hyperref[fig:7-1]{7-1} shows two bases of the same module. Observe that the parallelograms have equal area.

We note here that the unimodular matrices, or the corresponding linear transformations, form a group, the modular group\index{Modular group}.


% Original PDF physical page 284
\index{Unimodular transformation|)}\sourcepage{284}
\bookfigure{284}{1}{FIG. 7-1. Period module.}
\setcounter{section}{2}\setcounter{subsection}{2}\subsection{The Canonical Basis}\label{sec:7.2.3}
Among all possible bases of M it is possible to single out one, almost uniquely, to be called the canonical basis\index{Canonical basis|(}\index{Canonical basis}. It will not always be necessary, or even desirable, to use such a special basis, but it is important to know that one exists. Except for minor adjustments it will be the basis introduced in the course of the proof of Theorem~\hyperref[thm:7.1]{1}.

\begin{originalstatement}{Theorem}{2}\label{thm:7.2}
There exists a basis $(\omega_{1},\omega_{2})$ such that the ratio $\tau = \omega_{2} / \omega_{1}$ satisfies the following conditions: (i) Im $\tau >0$, (ii) $-\frac{1}{2} < \operatorname{Re}\tau \leq \frac{1}{2}$, (iii) $|\tau |\geq 1$, (iv) $\operatorname{Re}\tau \geq 0$ if $|\tau | = 1$. The ratio $\tau$ is uniquely determined by these conditions, and there is a choice of two, four, or six corresponding bases.
\end{originalstatement}
Proof. If we select $\omega_{1}$ and $\omega_{2}$ as in the proof of Theorem~\hyperref[thm:7.1]{1}, then $|\omega_{1}| \leq |\omega_{2}|$, $|\omega_{2}| \leq |\omega_{1} + \omega_{2}|$, and $|\omega_{2}| \leq |\omega_{1} - \omega_{2}|$. In terms of $\tau$ these conditions are equivalent to $|\tau| \geq 1$ and $|\operatorname{Re} \tau| \leq \frac{1}{2}$. If $\operatorname{Im} \tau < 0$ we replace $(\omega_{1}, \omega_{2})$ by $(- \omega_{1}, \omega_{2})$; this makes $\operatorname{Im} \tau > 0$ without changing the condition on $\operatorname{Re} \tau$. If $\operatorname{Re} \tau = -\frac{1}{2}$ we replace the basis by $(\omega_{1}, \omega_{1} + \omega_{2})$, and if $|\tau| = 1$, $\operatorname{Re} \tau < 0$ we replace it by $(- \omega_{2}, \omega_{1})$. After these minor changes all the conditions are satisfied.

Geometrically, the conditions (i) to (iv) mean that the point $\tau$ lies in the part of the complex plane shown in Fig.~\hyperref[fig:7-2]{7-2}. It is bounded by the circle $|\tau| = 1$ and the vertical lines $Re\tau = \pm\frac{1}{2}$, but only part of the
% Original PDF physical page 285
\sourcepage{285}
boundary is included. Although the set is not open, it is referred to as the fundamental region of the unimodular group.

We have seen that the most general change of basis is by a unimodular transformation. If the new ratio is $\tau'$ we obtain
\begin{equation*}
\tau^ {\prime} = \frac {a \tau + b}{c \tau + d}\tag{6}\label{eq:7.6}
\end{equation*}
with $ad - bc = \pm 1$. Simple computation gives
\begin{equation*}
\operatorname{Im} \tau^ {\prime} = \frac {\pm \operatorname{Im} \tau}{| c \tau + d | ^ {2}}\tag{7}\label{eq:7.7}
\end{equation*}
where the sign is the same as that of $ad - bc$.

Suppose that both $\tau$ and $\tau'$ are in the fundamental region. We shall show that they must then be equal. Our first remark is that it is the upper sign that is valid in \hyperref[eq:7.7]{(7)}, and hence $ad - bc = 1$. Second, because $\tau$ and $\tau'$ play symmetric roles, we are free to assume that $\operatorname{Im} \tau' \geq \operatorname{Im} \tau$. It then follows from \hyperref[eq:7.7]{(7)} that $|c\tau + d| \leq 1$. Because $c$ and $d$ are integers, there are very few possibilities for this inequality to hold.

One such possibility is to have $c = 0$, $d = \pm 1$. The relation $ad - bc = 1$ reduces to $ad = 1$, and because $a$ and $d$ are integers either $a = d = 1$ or $a = d = -1$. Equation \hyperref[eq:7.6]{(6)} becomes $\tau' = \tau \pm b$, and by condition (ii) it follows that $|b| = |\operatorname{Re} \tau' - \operatorname{Re} \tau| < 1$. Therefore, and because $b$ is an integer, $b = 0$ and $\tau' = \tau$.

Assume now that $c \neq 0$. The condition $|\tau + d/c| \leq 1 / |c|$ implies $|c| = 1$, for if $|c|$ were $\geq 2$, the point $\tau$ would be at a distance $\leq \frac{1}{2}$ from the real axis, which is obviously impossible, the nearest point in the fundamental region being at a distance $\sqrt{3}/2$. Thus $|\tau \pm d| \leq 1$, and a glance at Fig.~\hyperref[fig:7-2]{7-2} shows that this can occur only if $d = 0$ or $d = \pm 1$. The inequality $|\tau + 1| \leq 1$ is never fulfilled, for the point $e^{2\pi i/3}$ is not in the fundamental region, and $|\tau - 1| \leq 1$ only when $\tau = e^{\pi i/3}$. In the latter case $|c\tau + d| = 1$, and it follows from \hyperref[eq:7.7]{(7)} that $\operatorname{Im} \tau' = \operatorname{Im} \tau$ and hence, by the shape of the fundamental region, $\tau' = \tau$.

There remains only the case $d = 0$, $|c| = 1$. The condition $|\tau| \leq 1$ together with (iii) shows that $|\tau| = 1$. From $bc = -1$, it follows that $b/c = -1$ and $\tau' = \pm a - 1/\tau = \pm a - \bar{\tau}$. Hence $\operatorname{Re}\left(\tau + \tau'\right) = \pm a$, and by (ii) this is possible only for $a = 0$, in which case $\tau' = -1/\tau$. There is then a contradiction with (iv) unless $\tau = \tau' = i$.

We have proved that $\tau$ is unique. The canonical basis\index{Canonical basis} $(\omega_{1},\omega_{2})$ can always be replaced by $(-\omega_{1},-\omega_{2})$. There are other bases with the same $\tau$ if and only if $\tau$ is a fixed point of a unimodular transformation \hyperref[eq:7.6]{(6)}. This happens only for $\tau=i$ and $\tau=e^{\pi i/3}$; the former is a fixed point of $-1/\tau$, the latter of $-(\tau+1)/\tau$ and of $-1/(\tau+1)$. These are the multiple choices referred to in the theorem.


% Original PDF physical page 286
\index{Canonical basis|)}\sourcepage{286}
\bookfigure{286}{1}{FIG. 7-2. $\tau$-plane.}
\setcounter{section}{2}\setcounter{subsection}{3}\subsection{General Properties of Elliptic Functions}\label{sec:7.2.4}
In the following $f(z)$ will denote a meromorphic function which admits all numbers in the module M with basis $(\omega_{1},\omega_{2})$ as periods. We shall not assume that the basis is canonical, and it will not be required that M comprise all the periods.

It is convenient to say that $z_{1}$ is congruent to $z_{2}$, $z_{1} \equiv z_{2} \pmod{M}$, if the difference $z_{1} - z_{2}$ belongs to M, i.e., $z_{1} = z_{2} + n_{1}\omega_{1} + n_{2}\omega_{2}$. The function f takes the same values at congruent points, and may thus be regarded as a function on the congruence classes. A concrete way to make use of this property is to restrict the function to a parallelogram $P_{a}$ with vertices $a, a + \omega_{1}, a + \omega_{2}, a + \omega_{1} + \omega_{2}$. By including part of the boundary we may represent each congruence class by exactly one point in $P_{a}$, and then f is completely determined by its values on $P_{a}$. The choice of a is irrelevant, and we leave it free in order to attain, for instance, that f has no poles on the boundary of $P_{a}$.

\begin{originalstatement}{Theorem}{3}\label{thm:7.3}
An elliptic function\index{Elliptic function} without poles is a constant.
\end{originalstatement}
If $f(z)$ has no poles, it is bounded on the closure of $P_{a}$, and hence in the whole plane. By Liouville's theorem (Chap.~\hyperref[chap:4]{4}, Sec.~\hyperref[sec:4.2.3]{2.3}) it must reduce to a constant.

Because the poles have no accumulation point there are only finitely many poles in $P_{a}$. When we speak of the poles of an elliptic function we mean a full set of mutually incongruent poles. Multiplicities are counted in the usual manner.


% Original PDF physical page 287
\sourcepage{287}
\begin{originalstatement}{Theorem}{4}\label{thm:7.4}
The sum of the residues of an elliptic function\index{Elliptic function} is zero.
\end{originalstatement}
We may choose a so that none of the poles fall on the boundary of $P_{a}$. If the boundary $\partial P_{a}$ is traced in the positive sense, the sum of the residues at the poles in $P_{a}$ equals
\begin{equation*}
\frac {1}{2 \pi i} \int_ {\partial P _ {a}} f (z) d z.
\end{equation*}
Because f has periods $\omega_{1}$, $\omega_{2}$ the integral vanishes, for the integrals over opposite sides of the parallelogram cancel against each other.

As a consequence of the theorem there does not exist an elliptic function with a single simple pole.

\begin{originalstatement}{Theorem}{5}\label{thm:7.5}
A nonconstant elliptic function has equally many poles as it has zeros.
\end{originalstatement}
The poles and zeros of f are simple poles of $f'/f$, which is itself an elliptic function. The multiplicities are the residues of $f'/f$, counted positive for zeros and negative for poles. The theorem now follows from Theorem~\hyperref[thm:7.4]{4}.

If c is any constant, $f(z) - c$ has the same poles as $f(z)$. Therefore, all values are assumed equally many times. The number of incongruent roots of the equations $f(z) = c$ is called the order of the elliptic function.

\begin{originalstatement}{Theorem}{6}\label{thm:7.6}
The zeros $a_{1}, \ldots, a_{n}$ and poles $b_{1}, \ldots, b_{n}$ of an elliptic function satisfy $a_{1} + \cdots + a_{n} \equiv b_{1} + \cdots + b_{n} \pmod{M}$.
\end{originalstatement}
This is proved by considering the integral
\begin{equation*}
\frac {1}{2 \pi i} \int_ {\partial P _ {a}} \frac {z f ^ {\prime} (z)}{f (z)} d z\tag{8}\label{eq:7.8}
\end{equation*}
where we may again assume that there are no zeros or poles on the boundary. By the calculus of residues the integral equals $a_{1} + \cdots + a_{n} - b_{1} - \cdots - b_{n}$ provided that we choose the representative zeros and poles inside $P_{a}$. Consider the sides from a to $a + \omega_{1}$ and from $a + \omega_{2}$ to $a + \omega_{1} + \omega_{2}$. The corresponding part of the integral may be written
\begin{equation*}
\frac {1}{2 \pi i} \left(\int_ {a} ^ {a + \omega_ {1}} - \int_ {a + \omega_ {2}} ^ {a + \omega_ {1} + \omega_ {2}}\right) \frac {z f ^ {\prime} (z)}{f (z)} d z = - \frac {\omega_ {2}}{2 \pi i} \int_ {a} ^ {a + \omega_ {1}} \frac {f ^ {\prime} (z)}{f (z)} d z.
\end{equation*}
Except for the factor $-\omega_{2}$ the right-hand member represents the winding number around the origin of the closed curve described by $f(z)$ when z
% Original PDF physical page 288
\sourcepage{288}
varies\index{Weierstrass p-function@Weierstrass $\wp$-function|(}\index{p-function@$\wp$-function|(} from a to $a + \omega_{1}$. It is consequently an integer. The same applies to the other pair of opposite sides. Therefore the value of \hyperref[eq:7.8]{(8)} is of the form $n_{1}\omega_{1} + n_{2}\omega_{2}$, and the theorem is proved.

\setcounter{section}{2}\section{The Weierstrass Theory}\label{sec:7.3}
The simplest elliptic functions\index{Elliptic function} are of order 2, and such functions have either a double pole with residue zero, or two simple poles with opposite residues. We shall follow the classical example of Weierstrass, who chose a function with a double pole as the starting point of a systematic theory.

\setcounter{section}{3}\setcounter{subsection}{0}\subsection{\texorpdfstring{The Weierstrass $\wp$-function}{The Weierstrass ℘-function}}\label{sec:7.3.1}
We may as well place the pole at the origin, and since multiplication with a constant factor is clearly irrelevant, we may require that the singular part is $z^{-2}$. If $f$ is elliptic and has only this singularity at the origin and its congruent points, it is easy to see that $f$ must be an even function. Indeed, $f(z) - f(-z)$ has the same periods and no singularity. Therefore it must reduce to a constant, and on setting $z = \omega_1 / 2$ we conclude that the constant is zero.

A constant can be added at will, and we can therefore choose the constant term in the Laurent development about the origin to be zero. With this additional normalization $f(z)$ is uniquely determined, and it is traditionally denoted by a special typographical symbol $\wp(z)$. The Laurent development has the form
\begin{equation*}
\wp (z) = z ^ {- 2} + a _ {1} z ^ {2} + a _ {2} z ^ {4} + \dots .
\end{equation*}
So far all this is hypothetical, for we have not yet shown the existence of an elliptic function with this development. We shall follow the usual procedure in such cases, namely to postulate the existence and derive an explicit expression. The clue is to develop in partial fractions by the method in Chap.~\hyperref[chap:5]{5}, Sec.~\hyperref[sec:5.2]{2}. Our aim is to prove the formula
\begin{equation*}
\wp (z) = \frac {1}{z ^ {2}} + \sum_ {\omega \neq 0} \left(\frac {1}{(z - \omega) ^ {2}} - \frac {1}{\omega^ {2}}\right)\tag{9}\label{eq:7.9}
\end{equation*}
where the sum ranges over all $\omega = n_{1}\omega_{1} + n_{2}\omega_{2}$ except 0. Observe that $(z - \omega)^{-2}$ is the singular part at $\omega$, and that we have subtracted $\omega^{-2}$ in order to produce convergence.

Our first task is to verify that the series converges. If $|\omega| > 2|z|$, say, an immediate estimate gives
\begin{equation*}
\left| \frac {1}{(z - \omega) ^ {2}} - \frac {1}{\omega^ {2}} \right| = \left| \frac {z (2 \omega - z)}{\omega^ {2} (z - \omega) ^ {2}} \right| \leq \frac {10 | z |}{| \omega | ^ {3}}.
\end{equation*}
% Original PDF physical page 289
\sourcepage{289}
Therefore the series \hyperref[eq:7.9]{(9)} converges, uniformly on every compact set, provided that
\begin{equation*}
\sum_ {\omega \neq 0} \frac {1}{| \omega | ^ {3}} <   \infty .
\end{equation*}
This is indeed the case. Because $\omega_{2}/\omega_{1}$ is nonreal, there exists a k>0 such that $|n_{1}\omega_{1}+n_{2}\omega_{2}|\geq k(|n_{1}|+|n_{2}|)$ for all real pairs $(n_{1},n_{2})$. If we consider only integers there are 4n pairs $(n_{1},n_{2})$ with $|n_{1}|+|n_{2}|=n$. This gives
\begin{equation*}
\sum_ {\omega \neq 0} | \omega | ^ {- 3} \leq 4 k ^ {- 3} \sum_ {1} ^ {\infty} n ^ {- 2} <   \infty .
\end{equation*}
The next step is to prove that the right-hand side of \hyperref[eq:7.9]{(9)} has periods $\omega_{1}$ and $\omega_{2}$. Direct verification is relatively cumbersome. Instead we write, temporarily,
\begin{equation*}
f (z) = \frac {1}{z ^ {2}} + \sum_ {\omega \neq 0} \left(\frac {1}{(z - \omega) ^ {2}} - \frac {1}{\omega^ {2}}\right)\tag{10}\label{eq:7.10}
\end{equation*}
and obtain by termwise differentiation
\begin{equation*}
f ^ {\prime} (z) = - \frac {2}{z ^ {3}} - \sum_ {\omega \neq 0} \frac {2}{(z - \omega) ^ {3}} = - 2 \sum_ {\omega} \frac {1}{(z - \omega) ^ {3}}.
\end{equation*}
The last sum is obviously doubly periodic. Therefore $f(z + \omega_{1}) - f(z)$ and $f(z + \omega_{2}) - f(z)$ are constants. Because $f(z)$ is even (as seen from \hyperref[eq:7.10]{(10)}), it suffices to choose $z = -\omega_{1}/2$ and $z = -\omega_{2}/2$ to conclude that the constants are zero. We have thus proved that f has the asserted periods.

It follows now that $\wp(z) - f(z)$ is a constant, and by the form of the development at the origin the constant is zero. We have thereby proved the existence of $\wp(z)$, and also that it can be represented by the series \hyperref[eq:7.9]{(9)}. For convenient reference we display the important formula
\begin{equation*}
\wp^ {\prime} (z) = - 2 \sum_ {\omega} \frac {1}{(z - \omega) ^ {3}}.\tag{11}\label{eq:7.11}
\end{equation*}
\setcounter{section}{3}\setcounter{subsection}{1}\subsection{\texorpdfstring{The Functions $\zeta(z)$ and $\sigma(z)$}{The Functions ζ(z) and σ(z)}}\label{sec:7.3.2}
Because $\wp(z)$ has zero residues, it is the derivative of a single-valued function. It is traditional to denote the antiderivative of $\wp(z)$ by $-\zeta(z)$, and to normalize it so that it is odd. By use of \hyperref[eq:7.9]{(9)} we are led to the explicit expression
\begin{equation*}
\zeta (z) = \frac {1}{z} + \sum_ {\omega \neq 0} \left(\frac {1}{z - \omega} + \frac {1}{\omega} + \frac {z}{\omega^ {2}}\right).\tag{12}\label{eq:7.12}
\end{equation*}
% Original PDF physical page 290
\sourcepage{290}
The\index{Legendre relation|(} convergence is obvious, for apart from the term 1/z we obtain the new series by integration from 0 to z along any path that does not pass through the poles.

It is clear that $\zeta(z)$ satisfies conditions $\zeta(z + \omega_1) = \zeta(z) + \eta_1$, $\zeta(z + \omega_2) = \zeta(z) + \eta_2$, where $\eta_1$ and $\eta_2$ are constants. They are connected with $\omega_1, \omega_2$ by a very simple relation. To derive it we choose any $a \neq 0$ and observe that
\begin{equation*}
\frac {1}{2 \pi i} \int_ {\partial P _ {a}} \zeta (z) d z = 1,
\end{equation*}
by the residue theorem. The integral is easy to evaluate by adding the contributions from opposite sides of the parallelogram, and we obtain the equation
\begin{equation*}
\eta_ {1} \omega_ {2} - \eta_ {2} \omega_ {1} = 2 \pi i,
\end{equation*}
known as Legendre's relation.

The integration can be carried one step further provided that we use an exponential to eliminate the multiple-valuedness. Just as easily we can verify directly that the product
\begin{equation*}
\sigma (z) = z \prod_ {\omega \neq 0} \left(1 - \frac {z}{\omega}\right) e ^ {z / \omega + \frac {1}{2} (z / \omega) ^ {2}}\tag{13}\label{eq:7.13}
\end{equation*}
converges and represents an entire function which satisfies
\begin{equation*}
\sigma^ {\prime} (z) / \sigma (z) = \zeta (z).
\end{equation*}
The formula \hyperref[eq:7.13]{(13)} is a canonical product representation of $\sigma(z)$.

How does $\sigma(z)$ change when $z$ is replaced by $z + \omega_1$ or $z + \omega_2$? From
\begin{equation*}
\frac {\sigma^ {\prime} (z + \omega_ {1})}{\sigma (z + \omega_ {1})} = \frac {\sigma^ {\prime} (z)}{\sigma (z)} + \eta_ {1}
\end{equation*}
it follows at once that
\begin{equation*}
\sigma (z + \omega_ {1}) = C _ {1} \sigma (z) e ^ {\eta_ {1} z}
\end{equation*}
with constant $C_{1}$. To determine the constant we observe that $\sigma(z)$ is an odd function. On setting $z = -\omega_{1}/2$ the value of $C_{1}$ can be determined, and we find that $\sigma(z)$ satisfies
\begin{equation*}
\begin{array}{l} \sigma (z + \omega_ {1}) = - \sigma (z) e ^ {\eta_ {1} (z + \omega_ {1} / 2)} \\ \sigma (z + \omega_ {2}) = - \sigma (z) e ^ {\eta_ {2} (z + \omega_ {2} / 2)}. \end{array}\tag{14}\label{eq:7.14}
\end{equation*}
\subsection*{EXERCISES}
\originalnumber{1}{exercise:7.3.2.1.1}
1. Show that any even elliptic function\index{Elliptic function} with periods $\omega_{1}$, $\omega_{2}$ can be expressed in the form
\begin{equation*}
C \prod_ {k = 1} ^ {n} \frac {\wp (z) - \wp (a _ {k})}{\wp (z) - \wp (b _ {k})} \quad (C = \text { const. })
\end{equation*}
% Original PDF physical page 291
\index{Legendre relation|)}\sourcepage{291}
provided that 0 is neither a zero nor a pole. What is the corresponding form if the function either vanishes or becomes infinite at the origin?

\originalnumber{2}{exercise:7.3.2.1.2}
2. Show that any elliptic function\index{Elliptic function} with periods $\omega_{1}$, $\omega_{2}$ can be written as
\begin{equation*}
C \prod_ {k = 1} ^ {n} \frac {\sigma (z - a _ {k})}{\sigma (z - b _ {k})} \quad (C = \text { const. }).
\end{equation*}
Hint: Use \hyperref[eq:7.14]{(14)} and Theorem~\hyperref[thm:7.6]{6}.

\setcounter{section}{3}\setcounter{subsection}{2}\subsection{The Differential Equation}\index{Differential equation|(}\index{Differential equation}\label{sec:7.3.3}
By use of formula \hyperref[eq:7.12]{(12)} it is easy to derive the Laurent expansion of $\zeta(z)$ about the origin, and differentiation will then yield the corresponding expansion of $\wp(z)$. We have first
\begin{equation*}
\frac {1}{z - \omega} + \frac {1}{\omega} + \frac {z}{\omega^ {2}} = - \frac {z ^ {2}}{\omega^ {3}} - \frac {z ^ {3}}{\omega^ {4}} - \dots
\end{equation*}
and when we sum over all periods we obtain
\begin{equation*}
\zeta (z) = \frac {1}{z} - \sum_ {k = 2} ^ {\infty} G _ {k} z ^ {2 k - 1}
\end{equation*}
where we have written
\begin{equation*}
G _ {k} = \sum_ {\omega \neq 0} \frac {1}{\omega^ {2 k}}.
\end{equation*}
Observe that the corresponding sums of odd powers of the periods are zero, as was to be expected since $\zeta$ is an odd function. Because
\begin{equation*}
\wp (z) = - \zeta^ {\prime} (z)
\end{equation*}
we obtain further
\begin{equation*}
\wp (z) = \frac {1}{z ^ {2}} + \sum_ {k = 2} ^ {\infty} (2 k - 1) G _ {k} z ^ {2 k - 2}.
\end{equation*}
In the following computation we write down only the significant terms, since it is understood that the omitted terms are of higher order:
\begin{equation*}
\wp (z) = \frac {1}{z ^ {2}} + 3 G _ {2} z ^ {2} + 5 G _ {3} z ^ {4} + \dots
\end{equation*}
\begin{equation*}
\wp^ {\prime} (z) = - \frac {2}{z ^ {3}} + 6 G _ {2} z + 20 G _ {3} z ^ {3} + \dots
\end{equation*}
\begin{equation*}
\wp^ {\prime} (z) ^ {2} = \frac {4}{z ^ {6}} - \frac {24 G _ {2}}{z ^ {2}} - 80 G _ {3} + \dots
\end{equation*}
\begin{equation*}
4 \wp (z) ^ {3} = \frac {4}{z ^ {6}} + \frac {36 G _ {2}}{z ^ {2}} + 60 G _ {3} + \dots
\end{equation*}
\begin{equation*}
60 G _ {2} \wp (z) = \frac {60 G _ {2}}{z ^ {2}} + 0 + \dots .
\end{equation*}
% Original PDF physical page 292
\sourcepage{292}
The last three lines yield
\begin{equation*}
\wp^ {\prime} (z) ^ {2} - 4 \wp (z) ^ {3} + 60 G _ {2} \wp (z) = - 140 G _ {3} + \dots .
\end{equation*}
Here the left-hand side is a doubly periodic function, and the right-hand side has no poles. We may therefore conclude that
\begin{equation*}
\wp^ {\prime} (z) ^ {2} = 4 \wp (z) ^ {3} - 60 G _ {2} \wp (z) - 140 G _ {3}.
\end{equation*}
It is customary to set $g_{2} = 60G_{2}, g_{3} = 140G_{3}$ so that the equation becomes
\begin{equation*}
\wp^ {\prime} (z) ^ {2} = 4 \wp (z) ^ {3} - g _ {2} \wp (z) - g _ {3}.\tag{15}\label{eq:7.15}
\end{equation*}
This is a first-order differential equation\index{Differential equation} for $w = \wp(z)$. It can be solved explicitly, namely, by the formula
\begin{equation*}
z = \int^ {w} \frac {d w}{\sqrt {4 w ^ {3} - g _ {2} w - g _ {3}}} + \mathrm{constant},
\end{equation*}
which shows that $\wp(z)$ is the inverse of an elliptic integral. More accurately, this connection is expressed by the identity
\begin{equation*}
z - z _ {0} = \int_ {\wp (z _ {0})} ^ {\wp (z)} \frac {d w}{\sqrt {4 w ^ {3} - g _ {2} w - g _ {3}}}
\end{equation*}
where the path of integration is the image under $\wp$ of a path from $z_0$ to $z$ that avoids the zeros and poles of $\wp'(z)$, and where the sign of the square root must be chosen so that it actually equals $\wp'(z)$.

We recall that we encountered the relationship between elliptic functions\index{Elliptic function} and elliptic integrals already in connection with the conformal mapping of rectangles and certain triangles (Chap.~\hyperref[chap:6]{6}, Sec.~\hyperref[sec:6.2]{2}).

\subsection*{*EXERCISES}
The Weierstrass functions satisfy numerous identities which are best dealt with in an exercise section. They can be proved either by comparing two elliptic functions with the same zeros and poles (when $\sigma$-functions are involved), or by comparing elliptic functions with the same singular parts (when only $\wp$- and $\zeta$-functions are involved). The following sequence of formulas is so arranged that we need to resort to this method only once.

\originalnumber{1}{exercise:7.3.3.2.1}
1.
\begin{equation*}
\wp (z) - \wp (u) = - \frac {\sigma (z - u) \sigma (z + u)}{\sigma (z) ^ {2} \sigma (u) ^ {2}}\tag{16}\label{eq:7.16}
\end{equation*}
(Use \hyperref[eq:7.14]{(14)} to show that the right-hand member is a periodic function of $z$. Find the multiplicative constant by comparing the Laurent developments.)


% Original PDF physical page 293
\sourcepage{293}
\originalnumber{2}{exercise:7.3.3.2.2}
2.
\begin{equation*}
\frac {\wp^ {\prime} (z)}{\wp (z) - \wp (u)} = \zeta (z - u) + \zeta (z + u) - 2 \zeta (z).\tag{17}\label{eq:7.17}
\end{equation*}
(Follows from \hyperref[eq:7.16]{(16)} by taking logarithmic derivatives.)

\originalnumber{3}{exercise:7.3.3.2.3}
3.
\begin{equation*}
\zeta (z + u) = \zeta (z) + \zeta (u) + \frac {1}{2} \frac {\wp^ {\prime} (z) - \wp^ {\prime} (u)}{\wp (z) - \wp (u)}.\tag{18}\label{eq:7.18}
\end{equation*}
(This is a symmetrized version of \hyperref[eq:7.17]{(17)}.)

\originalnumber{4}{exercise:7.3.3.2.4}
4. The addition theorem\index{Addition theorem} for the $\wp$-function:
\begin{equation*}
\wp (z + u) = - \wp (z) - \wp (u) + \frac {1}{4} \left(\frac {\wp^ {\prime} (z) - \wp^ {\prime} (u)}{\wp (z) - \wp (u)}\right) ^ {2}.\tag{19}\label{eq:7.19}
\end{equation*}
(Differentiation of \hyperref[eq:7.18]{(18)} leads to a formula which contains $\wp''(z)$. It can be eliminated by \hyperref[eq:7.15]{(15)} which gives $\wp'' = 6\wp^2 - \frac{1}{2} g_2$. Symmetrization yields \hyperref[eq:7.19]{(19)}. Observe that this is an algebraic addition theorem, for $\wp'(z)$ and $\wp'(u)$ can be expressed algebraically through $\wp(z)$ and $\wp(u)$.)

\originalnumber{5}{exercise:7.3.3.2.5}
5. Prove
\begin{equation*}
\wp (2 z) = \frac {1}{4} \left(\frac {\wp^ {\prime \prime} (z)}{\wp^ {\prime} (z)}\right) ^ {2} - 2 \wp (z).
\end{equation*}
\originalnumber{6}{exercise:7.3.3.2.6}
6. Prove $\wp'(z) = -\sigma(2z)/\sigma(z)^4$.

\originalnumber{7}{exercise:7.3.3.2.7}
7. Prove that
\begin{equation*}
\left| \begin{array}{c c c} \wp (z) & \wp^ {\prime} (z) & 1 \\ \wp (u) & \wp^ {\prime} (u) & 1 \\ \wp (u + z) & - \wp^ {\prime} (u + z) & 1 \end{array} \right| = 0.
\end{equation*}
\setcounter{section}{3}\setcounter{subsection}{3}\subsection{\texorpdfstring{The Modular Function $\lambda(\tau)$}{The Modular Function λ(τ)}}\label{sec:7.3.4}
The differential equation\index{Differential equation} \hyperref[eq:7.15]{(15)} can also be written as
\begin{equation*}
\wp^ {\prime} (z) ^ {2} = 4 (\wp (z) - e _ {1}) (\wp (z) - e _ {2}) (\wp (z) - e _ {3}),\tag{20}\label{eq:7.20}
\end{equation*}
where $e_1, e_2, e_3$ are the roots of the polynomial $4w^3 - g_2w - g_3$.

To find the values of the $e_{k}$ we determine the zeros of $\wp'(z)$. The symmetry and periodicity of $\wp(z)$ imply $\wp(\omega_{1}-z)=\wp(z)$. Hence $\wp'(\omega_{1}-z)=-\wp'(z)$, from which it follows that $\wp'(\omega_{1}/2)=0$. Similarly $\wp'(\omega_{2}/2)=0$, and also $\wp'((\omega_{1}+\omega_{2})/2)=0$. The numbers $\omega_{1}/2$, $\omega_{2}/2$ and $(\omega_{1}+\omega_{2})/2$ are mutually incongruent modulo the periods. Therefore they are precisely the three zeros of $\wp'$, which is of order 3, and all the zeros are simple. When we compare with \hyperref[eq:7.20]{(20)} it follows that we can set
\begin{equation*}
e _ {1} = \wp (\omega_ {1} / 2), \quad e _ {2} = \wp (\omega_ {2} / 2), \quad e _ {3} = \wp ((\omega_ {1} + \omega_ {2}) / 2).\tag{21}\label{eq:7.21}
\end{equation*}
% Original PDF physical page 294
\index{Differential equation|)}\index{Weierstrass p-function@Weierstrass $\wp$-function|)}\index{p-function@$\wp$-function|)}\sourcepage{294}
It follows, moreover, and this is very important, that \emph{these roots are all distinct}. Indeed, $\wp(z)$ assumes each value $e_k$ with multiplicity 2, and if two of them were equal that value would be assumed four times in contradiction with the fact that $\wp$ is of order 2.

If we substitute $z = \omega_{1}/2$, $\omega_{2}/2$ and $(\omega_{1} + \omega_{2})/2$ in the definition \hyperref[eq:7.9]{(9)} of $\wp(z)$ it is seen at once that the $e_{k}$ are homogeneous of order -2 in $\omega_{1}$, $\omega_{2}$ (in other words, if the periods are multiplied by t, then the $e_{k}$ are multiplied by $t^{-2}$ ). We conclude that the quantity
\begin{equation*}
\lambda (\tau) = \frac {e _ {3} - e _ {2}}{e _ {1} - e _ {2}}\tag{22}\label{eq:7.22}
\end{equation*}
depends only on the ratio $\tau = \omega_{2}/\omega_{1}$, as indicated by our notation. It is quite clear from \hyperref[eq:7.9]{(9)} that $\lambda(\tau)$ is the quotient of two analytic functions in the upper half plane Im $\tau > 0$. Because $e_{1} \neq e_{2}$ it is actually analytic, rather than meromorphic; because $e_{2} \neq e_{3}$ it is never equal to 0, and because $e_{1} \neq e_{3}$ it is never equal to 1.

We shall study the dependence on $\tau$ in greater detail. If the periods are subjected to the unimodular transformation
\begin{equation*}
\begin{array}{r} \omega_ {2} ^ {\prime} = a \omega_ {2} + b \omega_ {1} \\ \omega_ {1} ^ {\prime} = c \omega_ {2} + d \omega_ {1} \end{array}\tag{23}\label{eq:7.23}
\end{equation*}
then, first of all, the $\wp$-function does not change. Therefore, by looking at \hyperref[eq:7.20]{(20)}, the roots $e_k$ can at most be permuted. Let us see what actually happens. It is clear from \hyperref[eq:7.23]{(23)} that $\omega_1'/2 \equiv \omega_1/2$ and $\omega_2'/2 \equiv \omega_2/2$ if $a \equiv d \equiv 1 \pmod{2}$ and $b \equiv c \equiv 0 \pmod{2}$. Under this condition the $e_k$ do not change, and we have shown that
\begin{equation*}
\lambda \left(\frac {a \tau + b}{c \tau + d}\right) = \lambda (\tau) \quad \text {   for   } \quad \left( \begin{array}{c c} a & b \\ c & d \end{array} \right) \equiv \left( \begin{array}{c c} 1 & 0 \\ 0 & 1 \end{array} \right) (\mathrm{mod} 2).\tag{24}\label{eq:7.24}
\end{equation*}
The transformations which satisfy the congruence relation in \hyperref[eq:7.24]{(24)} form a subgroup of the modular group (cf. Sec.~\hyperref[sec:7.2.2]{2.2}), known as the congruence subgroup\index{Congruence subgroup} mod 2. Equation \hyperref[eq:7.24]{(24)} asserts that $\lambda(\tau)$ is invariant under this subgroup. Quite generally, when an analytic or meromorphic function is invariant under a group of linear transformations, we call it an automorphic function\index{Automorphic function}. More specifically, a function which is automorphic with respect to a subgroup of the modular group is called a modular function\index{Modular function} (or an elliptic modular function\index{Elliptic modular function}).

We still have to determine the behavior of $\lambda(\tau)$ under a modular transformation that does not belong to the congruence subgroup. It is sufficient to consider matrices congruent mod 2 to $\begin{pmatrix}1 & 1 \\ 0 & 1\end{pmatrix}$ and $\begin{pmatrix}0 & 1 \\ 1 & 0\end{pmatrix}$ respectively, for all other types can be composed from these. In the first
% Original PDF physical page 295
\sourcepage{295}
case we obtain $\omega_2'/2 = (\omega_1 + \omega_2)/2$ and $\omega_1'/2 = \omega_1/2$; this means that $e_2$ and $e_3$ are interchanged, while $e_1$ remains fixed, and hence $\lambda$ goes over into $(e_2 - e_3)/(e_1 - e_3) = \lambda/(\lambda - 1)$. In the second case $\omega_2'/2 = \omega_1/2$, $\omega_1'/2 = \omega_2/2$, so that $e_1$ and $e_2$ are interchanged, and $\lambda$ goes over into $1 - \lambda$. Sample transformations are $\tau \to \tau + 1$ and $\tau \to -1/\tau$. We find that $\lambda(\tau)$ satisfies the functional equations
\begin{equation*}
\lambda (\tau + 1) = \frac {\lambda (\tau)}{\lambda (\tau) - 1}, \quad \lambda \left(- \frac {1}{\tau}\right) = 1 - \lambda (\tau).\tag{25}\label{eq:7.25}
\end{equation*}
\setcounter{section}{3}\setcounter{subsection}{4}\subsection{\texorpdfstring{The Conformal Mapping by $\lambda(\tau)$}{The Conformal Mapping by λ(τ)}}\label{sec:7.3.5}
For convenience we shall henceforth use the normalization $\omega_{1}=1$, $\omega_{2}=\tau$. With this choice of periods we obtain from \hyperref[eq:7.9]{(9)} and \hyperref[eq:7.21]{(21)}
\begin{equation*}
e _ {3} - e _ {2} = \sum_ {m, n = - \infty} ^ {\infty} \left[ \frac {1}{(m - \frac {1}{2} + (n + \frac {1}{2}) \tau) ^ {2}} - \frac {1}{(m + (n - \frac {1}{2}) \tau) ^ {2}} \right]\tag{26}\label{eq:7.26}
\end{equation*}
\begin{equation*}
e _ {1} - e _ {2} = \sum_ {m, n = - \infty} ^ {\infty} \left[ \frac {1}{(m - \frac {1}{2} + n \tau) ^ {2}} - \frac {1}{(m + (n - \frac {1}{2}) \tau) ^ {2}} \right]
\end{equation*}
where the double series are absolutely convergent. Our first observation is that these quantities are real when $\tau$ is purely imaginary (this is also true of the individual $e_{k}$ ). Indeed, when we replace $\tau$ by $-\tau$ the sums remain the same, except for a rearrangement of the terms. We conclude that $\lambda(\tau)$ is real on the imaginary axis.

Because $\begin{pmatrix} 1 & 2 \\ 0 & 1 \end{pmatrix}$ is in the congruence subgroup mod 2 we have $\lambda(\tau + 2) = \lambda(\tau)$. In other words, $\lambda$ has period 2. As we have seen in Sec. 2, this means that $\lambda(\tau)$ can be expressed as a function of $e^{\pi i\tau}$. It would not be difficult to determine the Fourier development, but we shall be content to show that $\lambda(\tau) \to 0$ as $\operatorname{Im} \tau \to \infty$.

To evaluate \hyperref[eq:7.26]{(26)} we sum first with respect to $m$. This summation can be carried out explicitly by use of the formula
\begin{equation*}
\frac {\pi^ {2}}{\sin^ {2} \pi z} = \sum_ {- \infty} ^ {\infty} \frac {1}{(z - m) ^ {2}}
\end{equation*}
(Chap.~\hyperref[chap:5]{5}, Sec.~\hyperref[sec:5.2.1]{2.1}, \hyperref[eq:7.9]{(9)}). We obtain at once
\begin{equation*}
e _ {3} - e _ {2} = \pi^ {2} \sum_ {n = - \infty} ^ {\infty} \left(\frac {1}{\cos^ {2} \pi (n - \frac {1}{2}) \tau} - \frac {1}{\sin^ {2} \pi (n - \frac {1}{2}) \tau}\right)\tag{27}\label{eq:7.27}
\end{equation*}
\begin{equation*}
e _ {1} - e _ {2} = \pi^ {2} \sum_ {n = - \infty} ^ {\infty} \left(\frac {1}{\cos^ {2} \pi n \tau} - \frac {1}{\sin^ {2} \pi (n - \frac {1}{2}) \tau}\right).
\end{equation*}
% Original PDF physical page 296
\sourcepage{296}
The series are strongly convergent, both for $n \to +\infty$ and $n \to -\infty$, for $|\cos n\pi\tau|$ and $|\sin n\pi\tau|$ are comparable to $e^{|n|\pi\operatorname{Im}\tau}$; the convergence is uniform for $\operatorname{Im}\tau \geq \delta > 0$.

The limits can now be taken termwise, and we find that $e_{3} - e_{2} \rightarrow 0$, $e_{1} - e_{2} \rightarrow \pi^{2}$ (from the term n = 0). Hence $\lambda(\tau) \rightarrow 0$ as Im $\tau \rightarrow \infty$, uniformly with respect to the real part of $\tau$. It follows further by the second equation \hyperref[eq:7.25]{(25)} that $\lambda(\tau) \rightarrow 1$ when $\tau$ approaches 0 along the imaginary axis.

We need one more piece of information, namely the order to which $\lambda(\tau)$ vanishes together with $e^{\pi i\tau}$. From \hyperref[eq:7.27]{(27)} the leading terms in $e_3 - e_2$ are the ones corresponding to $n = 0$ and $n = 1$. The sum of these terms is
\begin{equation*}
2 \pi^ {2} \left[ \frac {4 e ^ {\pi i \tau}}{(1 + e ^ {\pi i \tau}) ^ {2}} + \frac {4 e ^ {\pi i \tau}}{(1 - e ^ {\pi i \tau}) ^ {2}} \right]
\end{equation*}
and we conclude that
\begin{equation*}
\lambda (\tau) e ^ {- \pi i \tau} \rightarrow 16\tag{28}\label{eq:7.28}
\end{equation*}
for $\operatorname{Im} \tau \to \infty$.

In Fig.~\hyperref[fig:7-3]{7-3} the region $\Omega$ is bounded by the imaginary axis, the line $\operatorname{Re} \tau = 1$, and the circle $|\tau - \frac{1}{2}| = \frac{1}{2}$. The transformation $\tau + 1$ maps the imaginary axis on $\operatorname{Re} \tau = 1$, and $1 - 1/\tau$ maps $\operatorname{Re} \tau = 1$ on $|\tau - \frac{1}{2}| = \frac{1}{2}$. Since $\lambda(\tau)$ is real on the imaginary axis, it follows by virtue of the relations \hyperref[eq:7.25]{(25)} that it is real on the whole boundary of $\Omega$. Furthermore, $\lambda(\tau) \to 1$ as $\tau$ tends to 0 and $\lambda(\tau) \to \infty$ as $\tau$ tends to 1 inside $\Omega$.

We apply the argument principle to determine the number of times $\lambda(\tau)$ takes a nonreal value $w_{0}$ in $\Omega$. Cut off the corners of $\Omega$ by means of a
\bookfigure{296}{1}{FIG. 7-3}
% Original PDF physical page 297
\sourcepage{297}
horizontal line segment $\operatorname{Im} \tau = t_0$ and its images under the transformations $-1/\tau$ and $1 - 1/\tau$ (these images are circles tangent to the real axis). For sufficiently large $t_0$ it is clear that $\lambda(\tau) \neq w_0$ in the portions that have been cut off. The circle near $\tau = 1$ is mapped by $\lambda(\tau)$ on a curve $\lambda = \lambda(1 - 1/\tau) = 1 - 1/\lambda(\tau)$; where $\tau = s + it_0$, $0 \leq s \leq 1$; in view of \hyperref[eq:7.28]{(28)} this is approximately a large semicircle in the upper half plane. It is now evident that the image of the contour of the truncated region $\Omega$ has winding number 1 about $w_0$ if $\operatorname{Im} w_0 > 0$, and winding number 0 if $\operatorname{Im} w_0 < 0$. As a result $\lambda(\tau)$ takes every value in the upper half plane exactly once in $\Omega$, and no value in the lower half plane. This is also sufficient to guarantee that $\lambda(\tau)$ is monotone on the boundary of $\Omega$. Indeed, if it were not, the derivative $\lambda'(\tau)$ would vanish at a boundary point, and it would be impossible for a full semicircular neighborhood of that boundary point to be mapped into the upper half plane.

\begin{originalstatement}{Theorem}{7}\label{thm:7.7}
The modular function $\lambda(\tau)$ effects a one-to-one conformal mapping of the region $\Omega$ onto the upper half plane. The mapping extends continuously to the boundary in such a way that $\tau = 0, 1, \infty$ correspond to $\lambda = 1, \infty, 0$.
\end{originalstatement}
By reflection the region $\Omega'$ that is symmetric to $\Omega$ with respect to the imaginary axis is mapped onto the lower half plane, and thus both regions together correspond to the whole plane, except for the points 0 and 1.

We shall also prove:

\begin{originalstatement}{Theorem}{8}\label{thm:7.8}
Every point $\tau$ in the upper half plane is equivalent under the congruence subgroup mod 2 to exactly one point in $\bar{\Omega} \cup \Omega'$.
\end{originalstatement}
We refer to Fig.~\hyperref[fig:7-4]{7-4}. The reader is asked to verify that the region $\Delta$ is mapped on the shaded regions in the figure by means of the linear transformations $\tau, -1/\tau, \tau - 1, 1/(1 - \tau), (\tau - 1)/\tau, \tau/(1 - \tau)$ which we shall denote by $S_1, S_2, \ldots, S_6$. The matrices of the inverse transformations $S_k^{-1}$ ($k = 1, \ldots, 6$) are in order
\begin{equation*}
\begin{gathered}
\begin{pmatrix}1&0\\0&1\end{pmatrix},\quad
\begin{pmatrix}0&-1\\1&0\end{pmatrix},\quad
\begin{pmatrix}1&1\\0&1\end{pmatrix},\\[3pt]
\begin{pmatrix}1&-1\\1&0\end{pmatrix},\quad
\begin{pmatrix}0&1\\-1&1\end{pmatrix},\quad
\begin{pmatrix}1&0\\1&1\end{pmatrix}.
\end{gathered}
\end{equation*}
One recognizes readily that these matrices form a complete set of mutually incongruent matrices in the sense that every unimodular matrix is congruent mod 2 to exactly one of them. Precisely the same can be shown for the transformations $S_k' (k = 1, \ldots, 6)$ which map $\Delta'$ on the unshaded regions in the figure (the task of writing them down is left to the reader).


% Original PDF physical page 298
\index{Elliptic function|)}\sourcepage{298}
Together\index{Fundamental region|(} the 12 images of $\bar{\Delta}$ and $\bar{\Delta}'$ cover the set $\bar{\Omega} \cup \bar{\Omega}'$ (closures should be taken with respect to the open half plane).

Let $\tau$ be any point in the upper half plane. The set $\bar{\Delta} \cup \bar{\Delta}'$ can be identified with the closure of the shaded region in Fig. 7-0. Therefore, according to Theorem~\hyperref[thm:7.2]{2} there exists a modular transformation $S$ such that $S_{\tau}$ lies in $\bar{\Delta} \cup \bar{\Delta}'$. Suppose first that $S_{\tau}$ is in $\bar{\Delta}$. We know that the matrix of $S$ is congruent mod 2 to the matrix of an $S_k^{-1}$. It follows that the matrix of $T = S_k S$ is congruent to the identity matrix; in other words, $T$ belongs to the congruence subgroup. Since $S_{\tau}$ lies in $\bar{\Delta}$ we know further that $T_{\tau} = S_k(S_{\tau})$ lies in $\bar{\Omega} \cup \bar{\Omega}'$. The same reasoning applies if $S_{\tau} \in \bar{\Delta}'$. Thus there is always a $T_{\tau}$ in $\bar{\Omega} \cup \bar{\Omega}'$, and a trivial consideration shows that it can be chosen in $\bar{\Omega} \cup \Omega'$.

The uniqueness follows readily from the fact that the $S_{k}$ as well as the $S_{k}^{\prime}$ are mutually incongruent. We shall leave it to the reader to work out the details.

\subsection*{*EXERCISE}
\originalnumber{}{exercise:7.3.5.3.unnumbered}
Show that the function
\begin{equation*}
J (\tau) = \frac {4}{27} \frac {(1 - \lambda + \lambda^ {2}) ^ {3}}{\lambda^ {2} (1 - \lambda) ^ {2}}
\end{equation*}
is automorphic with respect to the full modular group. Where does it take the values 0 and 1, and with what multiplicities? Show that
\begin{equation*}
J (\tau) = \frac {- 4 (e _ {1} e _ {2} + e _ {2} e _ {3} + e _ {3} e _ {1}) ^ {3}}{(e _ {1} - e _ {2}) ^ {2} (e _ {2} - e _ {3}) ^ {2} (e _ {3} - e _ {1}) ^ {2}}.
\end{equation*}
Show also that $J(\tau)$ maps the region $\Delta$ in Fig.~\hyperref[fig:7-4]{7-4} onto a half plane.

\bookfigure{298}{1}{FIG. 7-4. Fundamental region of $\lambda(\tau)$.}
\index{Fundamental region|)}